\chapter{复数；角的度量.}

我们在这里不再完整地重述复数理论或四元数理论的历史发展，这两门理论本质上属于代数的领域（见第 72 页及第 62 页）；但对于虚数的几何表示，我们要说上几句，它在许多方面都构成了数学史上决定性的进步。

复数与平面的点之间的双射对应，其第一个清楚的观念无可争议地应归之于 C. F. 高斯\(^{1}\)，而尤其应当归功于他的，则是他第一个懂得把这一思想应用于复数理论，并且清楚地预见到 19 世纪的分析学家们将从它得到的好处。在 17 世纪和 18 世纪的过程中，数学家们逐渐达到这样的确信：使二次方程得以求解的虚数，同样也使任意次数的代数方程得以求解。18 世纪期间发表了这一定理的大量证明尝试；但即使撇开那些仅仅依赖循环论证的尝试不谈，也没有一个尝试不是会招致严重异议的。高斯在对这些尝试作了细致的考察、对其中的漏洞作了严格的批评之后，在他的博士论文（写于 1797 年，发表于 1799 年）中提出要最终给出一个严格的证明；他接过达朗贝尔顺便透露的一个想法（在后者 1746 年发表的证明\(^{2}\) 中，达朗贝尔指出，平面上那些点 \((a,b)\)——\(a+ib\) 是多项式 \(P(x+iy)=X(x,y)+iY(x,y)\) 的根——乃是曲线 X=0 与 Y=0 的交点；而高斯通过对这两条曲线的定性研究，证明其中一条曲线的一段连续弧把由另一条曲线所界出的两个不同区域中的点连接起来，并由此断定这两条曲线相交（{[}124 a{]}，第 III 卷，第 3 页；另见{[}124 b{]}））：这个证明以其清晰和独创，相对于先前的各种尝试构成一个可观的进步，而且无疑是把纯拓扑学的论证应用于代数问题的最早例子之一。\(^{3}\)

在他的博士论文中，高斯并没有明确地定义平面的点与虚数之间的对应；至于后者，以及它们在两个世纪里引起的那些关于“存在”的问题，他甚至采取一种相当保留的态度，有意把他的一切论证都以只出现实数的形式呈现出来。但是，他的证明中思想的进程，若不预先假定平面的点与复数之间的一种完全自觉的认同，就会是完全不可理解的；而他同一时期在数论和椭圆函数理论（复数在其中同样出现）方面的研究，只会加强这个假设。虚数的几何概念对他来说已熟悉到什么程度，它在他手中又会引出什么结果，这在（直到我们时代才发表的）那些把复数应用于初等几何问题求解的笔记（{[}124 a{]}，第 IV 卷，第 396 页及第 VIII 卷，第 307 页）中清楚地显示出来。更为明确的还有 1811 年致贝塞尔的信（{[}124 a{]}，第 VIII 卷，第 90–91 页），他在其中勾画了复变量函数积分理论的要点：“正像”，他说，“人们可以用一条无界直线表示实数的整个域，人们同样可以用一张无界平面来直观表示（“sinnlich machen”）一切量、即实数和虚数的整个域，在这平面上，由横坐标 a 和纵坐标 b 确定的每一点，同时表示量 \(a + ib\)。x 的一个值向另一个值的连续移动因此是沿一条线进行的，从而可以用无穷多种方式来完成\ldots”

但直到 1831 年，高斯（联系着“高斯整数”\(a + ib\)（其中 \(a\) 和 \(b\) 是整数）的引入）才就这一点以如此清楚的方式公开表明他的想法（{[}124 a{]}，第 II 卷，Theoria Residuorum Biquadraticorum, Commentatio secunda，第 38 节，第 109 页，及 Anzeige，第 174 页及以下）。在此期间，虚数几何表示的思想被两位默默无闻的研究者各自独立地重新发现，他们都是业余数学家，或多或少靠自学成才，这一发现也是他们两人对科学的唯一贡献，而且他们两人与他们那个时代的科学环境都几乎没有接触。由于这个缘故，他们的工作面临着被完全遗忘的严重危险；就时间在先的那一位而言，事情恰恰就这样发生了：丹麦人 C. 韦塞尔的小册子于 1798 年问世，构思和编写都非常清楚，却直到一个世纪之后才被人从湮没中发掘出来；而第二位，瑞士人 J. 阿尔冈，也几乎遭遇了同样的不幸，他只是出于偶然，才在 1813 年亲眼看到自己七年前出版的著作被重新发掘出来。\(^{4}\) 这部著作在《热尔岗年鉴》（Annales de Gergonne）上激起了一场活跃的讨论，在 1820 年到 1830 年间，这个论题在法国和英国成了好几篇出版物（出自一些相当不知名的作者之手）的对象；但当时还缺少一位大人物的权威来结束这些争论，并把数学家们集结到新观点上来；必须等到本世纪中叶，虚数的几何表示才终于得到普遍采纳：在德国是靠上引的高斯著作，在英国是靠哈密顿和凯莱关于超复系统的工作，在法国最后是靠柯西的附议\(^{5}\)；而这一切仅仅比黎曼以一次辉煌的推广进一步扩大几何学在解析函数理论中的作用、并同时创建拓扑学早了几年。

用角在圆上截出的弧来度量角，这种做法与角的概念本身一样古老，而且已为巴比伦人所知，我们沿用至今的角的单位——度——正是得自他们；此外，对他们来说，这只是 0 与 360°之间的角的度量问题，这对他们已经足够，因为他们使用角，主要是为了确定天体在其视轨道上一些预先确定的点处的位置，并编制供科学场合和占星场合使用的数表。

在古典时代的希腊几何学家那里，角的定义（《几何原本》第 I 卷，定义 8 与 9）受到的限制还要更大，因为它只适用于小于两直角的角；而另一方面，由于他们的比的理论和度量的理论有赖于对被度量的诸量之任意大的倍数进行比较，角对他们来说便不可能是一种可度量的量，尽管在他们那里自然可以找到等角的概念、一个角大于或小于另一个角的概念，以及当两角之和不超过两直角时两角之和的概念。正如同分数加法的情形一样，角的度量在他们眼中因此必定只是一种没有任何科学价值的经验程序。这个观点在阿基米德论螺线的那篇令人赞叹的论文（{[}5 b{]}，第 II 卷，第 1–12 页）中得到了很好的说明：由于不能用动径与角成比例来定义螺线，他给出了螺线的一种运动学描述（定义 1，第 44 页；参看命题 12 的陈述，第 46 页），而正如这部著作其余部分所表明的，他成功地从中提取了假使他掌握角的度量的一般概念便会得到的一切。至于希腊的天文学家们，在这一点上如同在许多别的点上一样，他们看来满足于步其巴比伦前辈的后尘。

在这里，也如同在实数概念的演变中一样（参见第 150 页），希腊科学衰落时期严格性精神的松弛带来了“朴素的”观点的回归，这种观点在某些方面比僵硬的欧几里得观念更接近于我们自己的观念。于是，一位轻率的插补者在欧几里得的书中插入了那条著名的命题（《几何原本》第 VI 卷，命题 33）：“各角与它们在圆上截出的弧成比例”\(^{6}\)，而一位对这条命题的“证明”作评注的匿名注释者则毫不犹豫地引入了等于一个圆周的任意大倍数的弧以及与这些弧相对应的角，当然不带任何证成（{[}107{]}，第 V 卷，第 357 页）。但是，甚至 16 世纪的韦达，在他发现方程 \(\sin nx = \sin \alpha\) 有多个根、从而似乎已触及我们现代的角概念的时候，也只求出了对应于小于两直角的角的那些根（{[}319{]}，第 305–313 页）。只有在 17 世纪，这一观点才被决定性地超越；而在牛顿发现 \(\sin x\) 与 \(\cos x\) 的级数展开、从而为这些函数提供了对变量的一切值都有效的表达式之后，人们便在欧拉那里、联系着“虚”数的对数，找到了关于任意角的度量概念的精确观念（{[}108 a{]}，(1)，第 XVII 卷，第 220 页）。

当然，用圆上一段弧的长度来度量一个角的那种经典定义不仅直观，而且本质上是正确的；然而，为使它严格化，需要曲线长度的概念，也就是说需要积分学。从所涉及的结构来看，那是一个十分迂回的程序，而且可以只用拓扑群理论的方法而无须别的方法；实指数函数和复指数函数由此显得出自同一个本源，即刻画“单参数群”的那条定理。

